% © 2026 Ing. David Jaroš — CC BY-NC-ND 4.0
%
% This work is licensed under a Creative Commons Attribution-NonCommercial-NoDerivatives
% 4.0 International License (CC BY-NC-ND 4.0).
% See LICENSE.md for full license text.
%
% File: docs/notes/symmetry_from_automorphisms.tex
% Purpose: Expository note — WHY UBT derives gauge symmetries from automorphisms
%          of ℂ⊗ℍ rather than postulating them.
%          Target reader: physics/engineering graduate (Ing. level, ČVUT).
% Status: expository / [STD] + references — NO new claims, NO new derivations.
% Author: Ing. David Jaroš
% Date: 2026-07-13
% Version: 1.0

\ifdefined\INCLUDEMODE
\else
\documentclass[12pt,a4paper]{article}
\usepackage[a4paper,margin=2.5cm]{geometry}
\usepackage{amsmath,amssymb,amsthm}
\usepackage{mathtools}
\usepackage{hyperref}
\usepackage{xcolor}
\usepackage{mdframed}
\usepackage{booktabs}
\usepackage{microtype}
\usepackage{array}

\hypersetup{
  colorlinks=true,
  linkcolor=blue!60!black,
  citecolor=green!50!black,
  urlcolor=blue!60!black
}

\newtheorem{theorem}{Theorem}[section]
\newtheorem{lemma}[theorem]{Lemma}
\newtheorem{proposition}[theorem]{Proposition}
\theoremstyle{definition}
\newtheorem{definition}[theorem]{Definition}
\theoremstyle{remark}
\newtheorem{remark}[theorem]{Remark}
\newtheorem{example}[theorem]{Example}

\newmdenv[backgroundcolor=yellow!15, linecolor=orange!70!black, linewidth=1.5pt]{keybox}
\newmdenv[backgroundcolor=blue!7, linecolor=blue!50, linewidth=1.5pt]{pointerbox}

% ---- document metadata ----
\title{%
  Why Gauge Symmetry Is a \emph{Theorem} in UBT:\\[4pt]
  \large Automorphisms of $\mathbb{C}\otimes\mathbb{H}$
  as the Origin of Gauge Structure}
\author{Ing.\ David Jaro\v{s}}
\date{July 2026 — v1.0 (expository note, \textbf{[STD]}+refs)}

\begin{document}
\maketitle
\tableofcontents
\bigskip

\fi % end INCLUDEMODE guard

% ===========================================================================
\section{Motivation: Einstein's Dream and the Missing Symmetry Principle}
\label{sec:motivation}
% ===========================================================================

\textit{This section frames the central question: what makes two fields
genuinely unified rather than notationally combined?}

\subsection{Einstein's 1945 attempt}

Albert Einstein spent the last decade of his life searching for a
\emph{unified field theory} that would merge gravitation and electromagnetism
into a single geometric structure.  His 1945 proposal~\cite{einstein1945}
replaced the symmetric metric tensor $g_{\mu\nu}$ of General Relativity (GR)
with a \emph{non-symmetric} (Hermitian) complex tensor
\begin{equation}
  \label{eq:einstein_g}
  g_{\mu\nu} = G_{\mu\nu} + i\,B_{\mu\nu},
\end{equation}
where $G_{\mu\nu}$ is the ordinary symmetric metric and $B_{\mu\nu}$ is an
antisymmetric real part intended to carry the electromagnetic field.
The strategy was seductive: if you can write gravity and electromagnetism as
real and imaginary parts of a single object, perhaps they are aspects of the
same thing.

\subsection{Why it failed}

The programme encountered three decisive obstacles, documented in the
post-humous analysis of Damour, Deser, and McCarthy~\cite{damour_deser_mccarthy1992}:
\begin{enumerate}
  \item \textbf{$B_{\mu\nu}$ is not the electromagnetic field.}
    The antisymmetric part satisfies equations that do not reduce to Maxwell's
    theory in any transparent limit.  Einstein himself acknowledged this.
  \item \textbf{No unifying symmetry principle.}
    The split $G + iB$ is \emph{by hand}: there is no algebraic reason that
    the real and imaginary parts should mix under the same symmetry group.
    In Einstein's own words (paraphrased from \cite{einstein1945}): the
    theory lacks the ``higher symmetry'' that would force the two fields to
    transform together.
  \item \textbf{Nonlinear ghosts.}
    The field equations for the full non-symmetric $g_{\mu\nu}$ contain
    additional degrees of freedom with uncontrolled kinetics~\cite{moffat_ngt1995},
    threatening the physical interpretation.
\end{enumerate}

\subsection{The central question}

\begin{keybox}
  \textbf{Key question.}  What does it mean for two fields to be
  \emph{genuinely unified} rather than \emph{notationally glued together}?
  The answer this note develops: they are genuinely unified when they are
  \emph{coordinates of a single algebraic object}, and when the symmetry
  group acting on that object is determined by the object's
  \textbf{algebraic structure}~---~not chosen by the physicist.
\end{keybox}

Unified Biquaternion Theory (UBT) answers this challenge differently: instead
of starting with a metric and splitting it, it starts with an algebra
($\mathbb{C}\otimes\mathbb{H}$, defined in Section~\ref{sec:biquaternions}) and
asks what symmetries the algebra itself forces onto any field built from it.
The gauge groups of the Standard Model emerge as the automorphism group of the
algebra.  This note explains why.

% ===========================================================================
\section{Quaternions and Biquaternions from Scratch}
\label{sec:biquaternions}
% ===========================================================================

\textit{This section defines $\mathbb{H}$ and $\mathbb{C}\otimes\mathbb{H}$,
gives their explicit $2\times 2$ matrix representation, and works through a
multiplication example.}

\subsection{The quaternion algebra $\mathbb{H}$}

A \textbf{quaternion} is an expression of the form
\begin{equation}
  q = a + b\,\mathbf{i} + c\,\mathbf{j} + d\,\mathbf{k},
  \qquad a,b,c,d \in \mathbb{R},
\end{equation}
where $1, \mathbf{i}, \mathbf{j}, \mathbf{k}$ are four basis elements satisfying
\begin{equation}
  \label{eq:quat_rules}
  \mathbf{i}^2 = \mathbf{j}^2 = \mathbf{k}^2 = \mathbf{i}\mathbf{j}\mathbf{k} = -1.
\end{equation}
From these four relations the full multiplication table follows (all eight
non-trivial products):

\begin{center}
\begin{tabular}{c|cccc}
  $\times$ & $1$ & $\mathbf{i}$ & $\mathbf{j}$ & $\mathbf{k}$ \\
  \hline
  $1$           & $1$  & $\mathbf{i}$  & $\mathbf{j}$  & $\mathbf{k}$  \\
  $\mathbf{i}$  & $\mathbf{i}$  & $-1$ & $\mathbf{k}$  & $-\mathbf{j}$ \\
  $\mathbf{j}$  & $\mathbf{j}$  & $-\mathbf{k}$ & $-1$ & $\mathbf{i}$  \\
  $\mathbf{k}$  & $\mathbf{k}$  & $\mathbf{j}$  & $-\mathbf{i}$ & $-1$ \\
\end{tabular}
\end{center}

The set of all quaternions with this multiplication is the
\textbf{quaternion algebra} $\mathbb{H}$ (after William~R.\ Hamilton, who
discovered it in 1843).  It is an associative, non-commutative algebra over
$\mathbb{R}$ of real dimension~4.

\medskip
\noindent\textit{Geometrically:} the ``imaginary'' part
$b\,\mathbf{i}+c\,\mathbf{j}+d\,\mathbf{k}$ behaves like a vector in
$\mathbb{R}^3$; the product of two pure imaginary quaternions gives the
cross product plus a scalar correction.  Rotations of $\mathbb{R}^3$ can be
written compactly as $\mathbf{v} \mapsto u\mathbf{v}u^{-1}$ with $|u|=1$
(see Section~\ref{sec:automorphisms}).

\subsection{The biquaternion algebra $\mathbb{C}\otimes\mathbb{H}$}

\textbf{Biquaternions} are obtained by allowing the coefficients to be complex:
\begin{equation}
  \label{eq:biquat}
  q = z_0 + z_1\,\mathbf{I} + z_2\,\mathbf{J} + z_3\,\mathbf{K},
  \qquad z_\mu \in \mathbb{C},
\end{equation}
where we write $\mathbf{I},\mathbf{J},\mathbf{K}$ for the quaternion basis when
complex coefficients are present, to avoid confusion with the ordinary complex
unit $i\in\mathbb{C}$.  The resulting algebra is
\begin{equation}
  \mathbb{C}\otimes_{\mathbb{R}}\mathbb{H}
  \;:=\;
  \bigl\{ z_0 + z_1\mathbf{I} + z_2\mathbf{J} + z_3\mathbf{K}
          \;\big|\; z_\mu\in\mathbb{C} \bigr\},
\end{equation}
a real algebra of dimension~8.  The complex unit $i$ commutes with
$\mathbf{I},\mathbf{J},\mathbf{K}$; the quaternion product rules
\eqref{eq:quat_rules} hold unchanged.

\subsection{The Pauli-matrix isomorphism}

\begin{proposition}[\textbf{[STD]}]
  \label{prop:pauli_iso}
  There is an algebra isomorphism
  $\mathbb{C}\otimes\mathbb{H} \cong \mathrm{M}_2(\mathbb{C})$
  (the algebra of $2\times 2$ complex matrices) given explicitly by
  \begin{equation}
    \label{eq:pauli_map}
    1 \;\mapsto\; I_2,
    \qquad
    \mathbf{I} \;\mapsto\; -i\sigma_1,
    \qquad
    \mathbf{J} \;\mapsto\; -i\sigma_2,
    \qquad
    \mathbf{K} \;\mapsto\; -i\sigma_3,
  \end{equation}
  where $I_2$ is the $2\times 2$ identity matrix and $\sigma_1,\sigma_2,\sigma_3$
  are the Pauli matrices
  \begin{equation}
    \label{eq:pauli_matrices}
    \sigma_1 = \begin{pmatrix}0&1\\1&0\end{pmatrix},
    \quad
    \sigma_2 = \begin{pmatrix}0&-i\\i&0\end{pmatrix},
    \quad
    \sigma_3 = \begin{pmatrix}1&0\\0&-1\end{pmatrix}.
  \end{equation}
\end{proposition}

\noindent
The isomorphism is a concrete translation between two languages for the same
algebra.  To verify it respects the quaternion rules, check the key product:
\[
  \mathbf{I}\mathbf{J} = (-i\sigma_1)(-i\sigma_2) = -\sigma_1\sigma_2 = -i\sigma_3 = \mathbf{K}.\;\checkmark
\]

\subsection{Worked multiplication example}

\begin{example}
  \label{ex:biquat_mult}
  Let $p = 1 + 2\mathbf{I}$ and $q = 3\mathbf{J} - i\mathbf{K}$
  (real and complex coefficients respectively).
  Compute $p\cdot q$ directly and via the matrix representation.

  \medskip
  \noindent\textbf{Direct:}
  \begin{align*}
    p\cdot q
    &= (1 + 2\mathbf{I})(3\mathbf{J} - i\mathbf{K}) \\
    &= 1\cdot3\mathbf{J} + 1\cdot(-i\mathbf{K})
       + 2\mathbf{I}\cdot3\mathbf{J} + 2\mathbf{I}\cdot(-i\mathbf{K}) \\
    &= 3\mathbf{J} - i\mathbf{K}
       + 6(\mathbf{I}\mathbf{J}) - 2i(\mathbf{I}\mathbf{K}).
  \end{align*}
  Using the table: $\mathbf{I}\mathbf{J}=\mathbf{K}$ and
  $\mathbf{I}\mathbf{K}=-\mathbf{J}$, so
  \[
    p\cdot q = 3\mathbf{J} - i\mathbf{K} + 6\mathbf{K} + 2i\mathbf{J}
             = (3+2i)\mathbf{J} + (6-i)\mathbf{K}.
  \]

  \noindent\textbf{Matrix check:}
  \[
    P = I_2 - 2i\sigma_1
      = \begin{pmatrix}1&-2i\\-2i&1\end{pmatrix},
    \quad
    Q = -3i\sigma_2 - i\cdot(-i\sigma_3)
      = -3i\begin{pmatrix}0&-i\\i&0\end{pmatrix}
        -\begin{pmatrix}1&0\\0&-1\end{pmatrix}
      = \begin{pmatrix}-1&-3\\3&1\end{pmatrix}.
      \]
  \[
    PQ = \begin{pmatrix}1&-2i\\-2i&1\end{pmatrix}
         \begin{pmatrix}-1&-3\\3&1\end{pmatrix}
       = \begin{pmatrix}-1-6i & -3-2i \\ 3+2i & 1+6i\end{pmatrix}.
  \]
  The result should map to $(3+2i)\mathbf{J}+(6-i)\mathbf{K}$, i.e.\
  $(3+2i)(-i\sigma_2) + (6-i)(-i\sigma_3)$:
  \begin{align*}
    (3+2i)(-i\sigma_2) + (6-i)(-i\sigma_3)
    &= (2-3i)\begin{pmatrix}0&-i\\i&0\end{pmatrix}
      +(-1-6i)\begin{pmatrix}1&0\\0&-1\end{pmatrix} \\
    &= \begin{pmatrix}0&-3-2i\\3+2i&0\end{pmatrix}
      +\begin{pmatrix}-1-6i&0\\0&1+6i\end{pmatrix} \\
    &= \begin{pmatrix}-1-6i & -3-2i \\ 3+2i & 1+6i\end{pmatrix}.
  \end{align*}
  Both computations agree. \hfill$\square$
\end{example}

The key message of this section: $\mathbb{C}\otimes\mathbb{H}$ is not an
abstract curiosity.  It is concretely the algebra of $2\times 2$ complex
matrices, already familiar from quantum mechanics via the Pauli matrices.

% ===========================================================================
\section{What an Automorphism Is}
\label{sec:automorphisms}
% ===========================================================================

\textit{This section defines algebra automorphisms, uses a crystal analogy to
show that the automorphism group is a \emph{property} of the algebra (not a
choice), and proves that $\mathrm{Aut}(\mathbb{H})\cong\mathrm{SO}(3)$.
It then explains that unit biquaternions generate $\mathrm{SL}(2,\mathbb{C})$,
the double cover of the Lorentz group, without any geometric input.}

\subsection{Definition}

\begin{definition}[Algebra automorphism]
  Let $\mathcal{A}$ be an algebra (a vector space with a bilinear product).
  A \textbf{automorphism} of $\mathcal{A}$ is a bijective linear map
  $\varphi:\mathcal{A}\to\mathcal{A}$ that preserves the product:
  \[
    \varphi(a\cdot b) = \varphi(a)\cdot\varphi(b)
    \quad\text{for all }a,b\in\mathcal{A}.
  \]
  The set of all automorphisms under composition forms a group,
  the \textbf{automorphism group} $\mathrm{Aut}(\mathcal{A})$.
\end{definition}

\subsection{The crystal analogy}

Imagine a crystal lattice of sodium chloride (table salt).
The lattice has a definite structure: every sodium ion is surrounded by six
chlorine ions at fixed angles and distances.
The \emph{symmetry group} of the crystal is the set of all rigid motions
(rotations, reflections) that map the lattice to itself.

Crucially: the symmetry group is not something a physicist \emph{chooses}
when studying the crystal.  It is a \emph{property} of the crystal that can
be \emph{discovered} by asking which transformations the lattice permits.
If you rotate NaCl by $90^{\circ}$ about the right axis, it looks the same; if you
rotate by $30^{\circ}$, it does not.

\begin{keybox}
  The automorphism group of an algebra plays exactly the same role:
  it is a \textbf{property} of the algebraic structure, determined by the
  multiplication table.  The physicist discovers it; they do not invent it.
\end{keybox}

\subsection{Automorphisms of $\mathbb{H}$: an explicit example}

For the quaternion algebra $\mathbb{H}$, every automorphism has a specific form.

\begin{theorem}[\textbf{[STD]} — see, e.g., \cite{adler1995}]
  \label{thm:aut_H}
  Every automorphism of $\mathbb{H}$ is an \emph{inner} automorphism:
  \[
    \varphi_u(q) = u\,q\,u^{-1},
    \qquad u\in\mathbb{H},\; u\neq 0.
  \]
  Two unit quaternions $u$ and $-u$ define the \emph{same} automorphism.
  Therefore
  \[
    \mathrm{Aut}(\mathbb{H})
    \;\cong\;
    \mathrm{Sp}(1)/\{\pm 1\}
    \;\cong\;
    \mathrm{SO}(3),
  \]
  where $\mathrm{Sp}(1)$ denotes the unit-sphere group of quaternions
  $\{u\in\mathbb{H} : |u|=1\}$.
\end{theorem}

\noindent\textbf{What does $u\,q\,u^{-1}$ do geometrically?}
Write $q = a + \mathbf{v}$ where $a\in\mathbb{R}$ is the scalar part and
$\mathbf{v} = b\mathbf{i}+c\mathbf{j}+d\mathbf{k}$ is the vector part.
Then $uqu^{-1} = a + u\mathbf{v}u^{-1}$ (the scalar is unchanged).  The
vector part $\mathbf{v}$ is \emph{rotated} in $\mathbb{R}^3$ by the rotation
corresponding to $u$.

\begin{example}[Rotation of the $(\mathbf{i},\mathbf{j},\mathbf{k})$ space]
  \label{ex:rotation}
  Take $u = \cos(\pi/4) + \sin(\pi/4)\,\mathbf{k}
          = \tfrac{1}{\sqrt{2}}(1 + \mathbf{k})$.
  This represents a $90^{\circ}$ rotation about the $\mathbf{k}$-axis.

  Check: $|u|^2 = \tfrac{1}{2}(1^2+1^2) = 1$, so $u^{-1}
  = \cos(\pi/4) - \sin(\pi/4)\mathbf{k}
  = \tfrac{1}{\sqrt{2}}(1-\mathbf{k})$.

  Apply $\varphi_u$ to $\mathbf{i}$:
  \begin{align*}
    u\,\mathbf{i}\,u^{-1}
    &= \tfrac{1}{\sqrt{2}}(1+\mathbf{k})\,\mathbf{i}\,
       \tfrac{1}{\sqrt{2}}(1-\mathbf{k}) \\
    &= \tfrac{1}{2}(1+\mathbf{k})\,(\mathbf{i} - \mathbf{i}\mathbf{k}).
  \end{align*}
  Using $\mathbf{i}\mathbf{k} = -\mathbf{j}$ (from the table):
  \[
    = \tfrac{1}{2}(1+\mathbf{k})(\mathbf{i}+\mathbf{j})
    = \tfrac{1}{2}(\mathbf{i}+\mathbf{j}+\mathbf{k}\mathbf{i}+\mathbf{k}\mathbf{j}).
  \]
  Using $\mathbf{k}\mathbf{i}=\mathbf{j}$ and $\mathbf{k}\mathbf{j}=-\mathbf{i}$:
  \[
    = \tfrac{1}{2}(\mathbf{i}+\mathbf{j}+\mathbf{j}-\mathbf{i})
    = \mathbf{j}.
  \]
  So $\mathbf{i}\mapsto\mathbf{j}$, confirming a $90^{\circ}$ rotation in the
  $(\mathbf{i},\mathbf{j})$ plane.  Similarly $\mathbf{j}\mapsto-\mathbf{i}$
  and $\mathbf{k}\mapsto\mathbf{k}$.\hfill$\square$
\end{example}

\subsection{Unit biquaternions and the Lorentz group}

Moving from $\mathbb{H}$ to $\mathbb{C}\otimes\mathbb{H}$ enlarges the group
significantly.  Let $\mathrm{SL}(2,\mathbb{C})$ denote the group of $2\times2$
complex matrices with determinant~1.

\begin{proposition}[\textbf{[STD]}]
  \label{prop:sl2c}
  Under the Pauli isomorphism \eqref{eq:pauli_map}, the unit-norm elements
  $\{u\in\mathbb{C}\otimes\mathbb{H} : \det(u)=1\}$ correspond precisely
  to $\mathrm{SL}(2,\mathbb{C})$.  Moreover,
  \[
    \mathrm{SL}(2,\mathbb{C}) \cong \widetilde{\mathrm{SO}}^+(1,3),
  \]
  the universal double cover of the proper orthochronous Lorentz group.
\end{proposition}

\begin{keybox}
  \textbf{The Lorentz group emerges from the algebra before any geometry is
  introduced.}  No metric, no spacetime manifold, no coordinate choice~---~only
  the multiplication table of $\mathbb{C}\otimes\mathbb{H}$.
\end{keybox}

This is the first instance of the algebra-first principle: symmetry is not
postulated; it is read off from the structure.

% ===========================================================================
\section{From Global to Local: Gauge Fields}
\label{sec:gauging}
% ===========================================================================

\textit{This section explains the gauging procedure in one page and emphasises
that WHICH group can be gauged is dictated by the algebra.}

\subsection{Global vs.\ local symmetry}

Suppose a field $\psi(x)$ is invariant under transformations
$\psi(x)\mapsto u\,\psi(x)$ for a fixed (``global'') group element $u$.
This means $u$ is the same at every spacetime point $x$.

Now suppose we want the symmetry to hold \emph{independently at each
point}: $u = u(x)$.  This is the step from global to \textbf{local symmetry}
(also called gauging).  The problem is that the ordinary derivative
$\partial_\mu\psi$ does not transform cleanly:
\begin{equation}
  \partial_\mu (u(x)\psi(x))
  = u(x)\,\partial_\mu\psi(x)
    + (\partial_\mu u(x))\,\psi(x).
\end{equation}
The unwanted extra term $(\partial_\mu u)\psi$ spoils covariance unless we
introduce a \textbf{compensating field} $A_\mu$, a connection, that transforms
as
\begin{equation}
  A_\mu \;\mapsto\; u\,A_\mu\,u^{-1} - (\partial_\mu u)\,u^{-1}.
\end{equation}
The \textbf{covariant derivative}
$D_\mu\psi := \partial_\mu\psi + A_\mu\psi$
then transforms as $D_\mu\psi\mapsto u\,D_\mu\psi$.

\subsection{Which group can be gauged?}

The connection field $A_\mu$ must take values in the Lie algebra of the group
being gauged.  The question ``which group?'' therefore reduces to:

\begin{center}
  \textit{What are the automorphisms of the field algebra?}
\end{center}

In UBT, the field $\Theta(q,\tau)$ takes values in $\mathbb{C}\otimes\mathbb{H}$.
The automorphisms of this algebra that preserve its involution structure
(the $\mathbb{Z}_2^3$ structure; see canonical file
\texttt{canonical/algebra/\allowbreak involutions\_Z2xZ2xZ2.tex}) are:
\begin{itemize}
  \item The $\mathrm{SO}(3)$ rotations of the quaternionic part
    (Theorem~\ref{thm:aut_H});
  \item Their complexification $\mathrm{SL}(2,\mathbb{C})$
    (Proposition~\ref{prop:sl2c});
  \item Additional sector automorphisms linked to the $\mathbb{Z}_2^3$
    involution structure that yield $\mathrm{SU}(3)\times\mathrm{SU}(2)_L$
    (proved in the canonical layer, Section~\ref{sec:pointers}).
\end{itemize}

\subsection{Contrast with Chamseddine's approach}

Chamseddine's unified theories (2001, 2006~\cite{chamseddine2001,chamseddine2006})
work with a group $\mathrm{U}(1,3)$ or a spectral-triple action functional.
The group structure is \emph{postulated} as input to the action principle;
it is then shown to contain Standard Model gauge bosons.

In UBT the logic is reversed:

\begin{keybox}
\sloppy
  The Standard Model gauge group
  $\mathrm{SU}(3)\times\mathrm{SU}(2)_L\times\mathrm{U}(1)_Y$
  is the automorphism group of the algebra.  It is \textbf{not} an input;
  it is a \textbf{theorem} about $\mathbb{C}\otimes\mathbb{H}$.
  (Proofs: \texttt{canonical/su3\_derivation/\allowbreak su3\_from\_involutions.tex},
  \texttt{canonical/chirality/\allowbreak step4\_no\_wr\_derivation.tex}.)
\end{keybox}

This is exactly the ``higher symmetry principle'' Einstein identified as
missing from his 1945 theory.

% ===========================================================================
\section{The Tetrad and Why Components Cannot Move Independently}
\label{sec:tetrad}
% ===========================================================================

\textit{This section defines the tetrad pedagogically, presents the UBT metric
construction, and explains the central argument: under an automorphism, the
components of the biquaternionic tetrad mix according to the fixed
multiplication table of the algebra — they cannot be varied independently.}

\subsection{What is a tetrad?}

In GR, the metric $g_{\mu\nu}(x)$ encodes distances and angles at each
spacetime point.  A \textbf{tetrad} (or \emph{vierbein}, German for
``four-leg'') is a set of four vector fields $e_\mu^{\;a}(x)$ that provide a
\emph{local inertial frame} at every point.  They are related to the metric by
\begin{equation}
  \label{eq:metric_tetrad}
  g_{\mu\nu} = e_\mu^{\;a}\,e_\nu^{\;b}\,\eta_{ab},
\end{equation}
where $\eta_{ab} = \mathrm{diag}(-1,+1,+1,+1)$ is the Minkowski metric.
The index $a\in\{0,1,2,3\}$ is a ``flat'' (local frame) index;
$\mu,\nu$ are curved spacetime indices.

\medskip
\noindent\textbf{Mnemonic}: the tetrad is a ``square root of the metric.''
Equation~\eqref{eq:metric_tetrad} is the matrix-product version of the identity
$\mathbf{e}\cdot\mathbf{e}^T = g$.

\medskip
\noindent\textbf{Why the tetrad is needed:} the Dirac equation for a spin-$\frac{1}{2}$
particle requires a local inertial frame at each point to define what
``spinor'' means.  Without the tetrad, there is no natural way to couple
fermions to curved spacetime.

\subsection{The biquaternionic tetrad in UBT}

In UBT the tetrad is promoted to a biquaternion-valued field:
\begin{equation}
  E_\mu(x) \in \mathbb{C}\otimes\mathbb{H},
  \qquad \mu = 0,1,2,3.
\end{equation}
The metric is reconstructed as (notation: $\mathrm{Sc}(\cdot)$ denotes the
scalar part of a biquaternion, $E^\dagger$ the biquaternionic conjugate):
\begin{equation}
  \label{eq:ubt_metric}
  \mathcal{G}_{\mu\nu} = \mathrm{Sc}(E_\mu\,E_\nu^\dagger).
\end{equation}
In the real limit this reduces to the standard GR metric.
(Proved in \texttt{canonical/gr\_closure/\allowbreak step1\_metric\_bridge.tex}.)

\begin{example}[Scalar extraction — $2\times 2$ matrix version]
  \label{ex:tetrad_scalar}
  Let $E_0 = 1 + \mathbf{I}$ (a simple quaternionic tetrad component; in
  matrix form $E_0 = I_2 - i\sigma_1 = \bigl(\begin{smallmatrix}1&-i\\-i&1\end{smallmatrix}\bigr)$).
  The biquaternionic conjugate is $E_0^\dagger = 1 - \mathbf{I}$
  ($= \bigl(\begin{smallmatrix}1&i\\i&1\end{smallmatrix}\bigr)$).
  Then
  \[
    E_0\,E_0^\dagger
    = (1+\mathbf{I})(1-\mathbf{I})
    = 1 - \mathbf{I}^2 = 1 - (-1) = 2.
  \]
  The scalar part is $\mathrm{Sc}(2) = 2$, and the metric component
  $\mathcal{G}_{00} = 2$.  In the matrix picture:
  \[
    \begin{pmatrix}1&-i\\-i&1\end{pmatrix}
    \begin{pmatrix}1&i\\i&1\end{pmatrix}
    = \begin{pmatrix}1+1&-i+i\\-i+i&1+1\end{pmatrix}
    = \begin{pmatrix}2&0\\0&2\end{pmatrix}
    = 2I_2.
  \]
  The scalar part (half the trace) is indeed $2$.  \hfill$\square$
\end{example}

\subsection{Why the components cannot move independently}

Now we reach the central argument.  Restrict to unitary
$u\in\mathrm{SU}(2)\subset\mathrm{SL}(2,\mathbb{C})$ (so $u^\dagger=u^{-1}$),
and consider the transformation
\begin{equation}
  E_\mu \;\mapsto\; u\,E_\mu\,u^\dagger,
  \qquad u^\dagger=u^{-1}.
\end{equation}
Under this transformation the metric becomes
\begin{align}
  \mathcal{G}'_{\mu\nu}
  &= \mathrm{Sc}\bigl((uE_\mu u^\dagger)(uE_\nu u^\dagger)^{\dagger}\bigr)
   = \mathrm{Sc}(u E_\mu E_\nu^\dagger u^\dagger).
\end{align}
Because $u$ is unitary and the scalar part satisfies
$\mathrm{Sc}(uXu^\dagger) = \mathrm{Sc}(uXu^{-1}) = \mathrm{Sc}(X)$ in this
unitary case, the metric is \emph{invariant}:
$\mathcal{G}'_{\mu\nu} = \mathcal{G}_{\mu\nu}$.  For general
$\mathrm{SL}(2,\mathbb{C})$ elements (boost sector), $\mathcal{G}_{\mu\nu}$
is not invariant but transforms covariantly as a tensor; see
\texttt{canonical/gr\_closure/step1\_metric\_bridge.tex} for the full
Lorentz-covariant treatment.

The crucial point is \emph{how} $u$ acts on $E_\mu$.  The tetrad component
can be decomposed as
\begin{equation}
  E_\mu = s_\mu + v_\mu^a\,\mathbf{e}_a
  \qquad
  (s_\mu\in\mathbb{C}\text{ scalar part},\;
   v_\mu^a\in\mathbb{C}\text{ vector part}),
\end{equation}
but the action $E_\mu\mapsto uE_\mu u^\dagger$ mixes $s_\mu$ and $v_\mu^a$
\emph{according to the quaternion multiplication table}~---~which is a fixed
algebraic law, not a choice.

In Einstein's 1945 theory the ``metric'' $G_{\mu\nu}+iB_{\mu\nu}$ has 16
independent real components.  Nothing in the theory prevents a physicist from
varying $G$ and $B$ independently; they are not coordinates of a single
algebraic object.  In UBT:

\begin{keybox}
  The components of $E_\mu$ are entries in a single element of
  $\mathbb{C}\otimes\mathbb{H}$.  Any symmetry transformation must respect the
  algebra's multiplication table.  The metric, the connection, and the gauge
  fields all arise from \emph{one} object under the \emph{one} automorphism
  group of that object.  This is the higher symmetry principle Einstein sought.
\end{keybox}

% ===========================================================================
\section{Pointers to Canonical Proofs}
\label{sec:pointers}
% ===========================================================================

\textit{This note is expository; it establishes no new results.  Every claim
above is either a standard mathematical theorem or refers to a proof in the
canonical layer.  The table below maps each main claim to its proof file.}

\begin{pointerbox}
\renewcommand{\arraystretch}{1.4}
\small
\begin{tabular}{>{\raggedright\arraybackslash}p{5.0cm}
                >{\raggedright\arraybackslash}p{1.8cm}
                >{\raggedright\arraybackslash}p{6.2cm}}
  \toprule
  \textbf{Claim} & \textbf{Level} & \textbf{Canonical source} \\
  \midrule
  Pauli isomorphism
    $\mathbb{C}\otimes\mathbb{H}\cong\mathrm{M}_2(\mathbb{C})$
  & {[STD]}
  & \texttt{canonical/algebra/}\linebreak\texttt{biquaternion\_algebra.tex} \\
  $\mathrm{Aut}(\mathbb{H})\cong\mathrm{SO}(3)$
  & {[STD]}
  & \texttt{canonical/algebra/}\linebreak\texttt{biquaternion\_algebra.tex};
    \cite{adler1995} §1 \\
  $\mathrm{SL}(2,\mathbb{C})$ is double cover of Lorentz group
  & {[STD]}
  & Standard QFT, e.g.\ \cite{mtw1973} \\
  $\mathbb{Z}_2^3$ involution structure of $\mathbb{C}\otimes\mathbb{H}$
  & {[L0]}
  & \texttt{canonical/algebra/}\linebreak\texttt{involutions\_Z2xZ2xZ2.tex} \\
  $\mathrm{SU}(3)$ from involutions
  & {[L1]}
  & \texttt{canonical/su3\_derivation/}\linebreak\texttt{su3\_from\_involutions.tex} \\
  Chirality: $\mathrm{SU}(2)_L$ selection, $\mathrm{SU}(2)_R$ exclusion
  & {[L1 cond.]}
  & \texttt{canonical/chirality/}\linebreak\texttt{step4\_no\_wr\_derivation.tex} \\
  Metric from biquaternionic tetrad
    $\mathcal{G}_{\mu\nu}=\mathrm{Sc}(E_\mu E_\nu^\dagger)$
  & {[L1]}
  & \texttt{canonical/gr\_closure/}\linebreak\texttt{step1\_metric\_bridge.tex} \\
  GR field equations from UBT in real limit
  & {[L1]}
  & \texttt{papers/UBT\_GR\_Submission.tex}~§3;
    \texttt{canonical/gr\_closure/}\linebreak\texttt{step3\_einstein\_with\_matter.tex} \\
  \bottomrule
\end{tabular}
\end{pointerbox}

\bigskip
\noindent
Readers wishing to check any claim in full detail should consult the canonical
layer files listed above.  The files compile standalone via
the \verb|\INCLUDEMODE| guard or as part of the main compilation.

% ===========================================================================
\section*{Note on Authorship and Method}
\label{sec:authorship}
\addcontentsline{toc}{section}{Note on Authorship and Method}
% ===========================================================================

The core ideas presented in this note~---~the biquaternionic field
formulation and the algebra-first approach to gauge symmetry~---~were
developed by the author in 2013--2015, predating this exposition.
See \texttt{docs/PRIORITY.md} and the research archive commit history for
the original development timeline.

The text of this exposition was drafted with AI assistance (GitHub Copilot)
under the author's direction and review.  All mathematical content is either
standard literature (labelled \textbf{[STD]}) or references proofs in the
canonical layer of this repository.

\begin{itemize}
  \item Historical context on non-symmetric unified field theory:
    Silberstein \cite{silberstein1912},
    Einstein \cite{einstein1945},
    Einstein--Straus \cite{einstein_straus1946},
    Damour--Deser--McCarthy \cite{damour_deser_mccarthy1992},
    Moffat \cite{moffat_ngt1995}.
  \item Algebra-first unification via noncommutative geometry:
    Chamseddine (2001) \cite{chamseddine2001},
    Chamseddine (2006) \cite{chamseddine2006}.
  \item Quaternionic quantum mechanics: \cite{adler1995}.
  \item General Relativity conventions: \cite{mtw1973,wald1984}.
\end{itemize}

% ===========================================================================
% Bibliography
% ===========================================================================

\ifdefined\INCLUDEMODE
\else

\begin{thebibliography}{99}

\bibitem{einstein1915}
A.~Einstein,
\textit{Die Feldgleichungen der Gravitation},
Sitzungsberichte der Preussischen Akademie der Wissenschaften,
pp.\ 844--847, 1915.

\bibitem{silberstein1912}
L.~Silberstein,
\textit{Quaternionic form of relativity},
Philosophical Magazine, vol.~23, pp.\ 790--809, 1912.

\bibitem{einstein1945}
A.~Einstein,
\textit{A Generalization of the Relativistic Theory of Gravitation},
Annals of Mathematics, vol.~46, no.~4, pp.\ 578--584, 1945.

\bibitem{einstein_straus1946}
A.~Einstein and E.~G.~Straus,
\textit{A Generalization of the Relativistic Theory of Gravitation, II},
Annals of Mathematics, vol.~47, no.~4, pp.\ 731--741, 1946.

\bibitem{damour_deser_mccarthy1992}
T.~Damour, S.~Deser, and J.~McCarthy,
\textit{Nonsymmetric gravity theories: Inconsistencies and a cure},
Physical Review D, vol.~45, p.~3289, 1992.

\bibitem{moffat_ngt1995}
J.~W.~Moffat,
\textit{Nonsymmetric gravitational theory},
Journal of Mathematical Physics, vol.~36, no.~7, pp.\ 3722--3732, 1995.

\bibitem{chamseddine2001}
A.~H.~Chamseddine,
\textit{Complexified Gravity in Noncommutative Spaces},
Communications in Mathematical Physics, vol.~218, no.~2, pp.\ 283--292, 2001.

\bibitem{chamseddine2006}
A.~H.~Chamseddine,
\textit{Gravity in Complex Hermitian Space-Time},
Annales Henri Poincar\'{e}, 2006.
\texttt{arXiv:hep-th/0610099}.

\bibitem{adler1995}
S.~L.~Adler,
\textit{Quaternionic Quantum Mechanics and Quantum Fields},
Oxford University Press, New York, 1995.

\bibitem{mtw1973}
C.~W.~Misner, K.~S.~Thorne, and J.~A.~Wheeler,
\textit{Gravitation},
W.~H.~Freeman, San Francisco, 1973.

\bibitem{wald1984}
R.~M.~Wald,
\textit{General Relativity},
University of Chicago Press, Chicago, 1984.

\end{thebibliography}

\end{document}
\fi % end INCLUDEMODE guard
